\documentclass{amsart}
% For dealing with references we use the comment environment
\usepackage{verbatim}
\newenvironment{reference}{\comment}{\endcomment}
%\newenvironment{reference}{}{}
\newenvironment{slogan}{\comment}{\endcomment}
\newenvironment{history}{\comment}{\endcomment}

% For commutative diagrams we use Xy-pic
\usepackage[all]{xy}

% We use 2cell for 2-commutative diagrams.
\xyoption{2cell}
\UseAllTwocells

% We use multicol for the list of chapters between chapters
\usepackage{multicol}

% This is generally recommended for better output
\usepackage{lmodern}
\usepackage[T1]{fontenc}

% For cross-file-references

% Package for hypertext links:
\usepackage{hyperref}

% For any local file, say "hello.tex" you want to link to please

% Theorem environments.
%
\theoremstyle{plain}
\newtheorem{theorem}[subsection]{Theorem}
\newtheorem{proposition}[subsection]{Proposition}
\newtheorem{lemma}[subsection]{Lemma}

\theoremstyle{definition}
\newtheorem{definition}[subsection]{Definition}
\newtheorem{example}[subsection]{Example}
\newtheorem{exercise}[subsection]{Exercise}
\newtheorem{situation}[subsection]{Situation}

\theoremstyle{remark}
\newtheorem{remark}[subsection]{Remark}
\newtheorem{remarks}[subsection]{Remarks}

\numberwithin{equation}{subsection}

% Macros
%
\def\lim{\mathop{\mathrm{lim}}\nolimits}
\def\colim{\mathop{\mathrm{colim}}\nolimits}
\def\Spec{\mathop{\mathrm{Spec}}}
\def\Hom{\mathop{\mathrm{Hom}}\nolimits}
\def\Ext{\mathop{\mathrm{Ext}}\nolimits}
\def\SheafHom{\mathop{\mathcal{H}\!\mathit{om}}\nolimits}
\def\SheafExt{\mathop{\mathcal{E}\!\mathit{xt}}\nolimits}
\def\Sch{\mathit{Sch}}
\def\Mor{\mathop{\mathrm{Mor}}\nolimits}
\def\Ob{\mathop{\mathrm{Ob}}\nolimits}
\def\Sh{\mathop{\mathit{Sh}}\nolimits}
\def\NL{\mathop{N\!L}\nolimits}
\def\CH{\mathop{\mathrm{CH}}\nolimits}
\def\proetale{{pro\text{-}\acute{e}tale}}
\def\etale{{\acute{e}tale}}
\def\QCoh{\mathit{QCoh}}
\def\Ker{\mathop{\mathrm{Ker}}}
\def\Im{\mathop{\mathrm{Im}}}
\def\Coker{\mathop{\mathrm{Coker}}}
\def\Coim{\mathop{\mathrm{Coim}}}

% Boxtimes
%
\DeclareMathSymbol{\boxtimes}{\mathbin}{AMSa}{"02}

%
% Macros for moduli stacks/spaces
%
\def\QCohstack{\mathcal{QC}\!\mathit{oh}}
\def\Cohstack{\mathcal{C}\!\mathit{oh}}
\def\Spacesstack{\mathcal{S}\!\mathit{paces}}
\def\Quotfunctor{\mathrm{Quot}}
\def\Hilbfunctor{\mathrm{Hilb}}
\def\Curvesstack{\mathcal{C}\!\mathit{urves}}
\def\Polarizedstack{\mathcal{P}\!\mathit{olarized}}
\def\Complexesstack{\mathcal{C}\!\mathit{omplexes}}
% \Pic is the operator that assigns to X its picard group, usage \Pic(X)
% \Picardstack_{X/B} denotes the Picard stack of X over B
% \Picardfunctor_{X/B} denotes the Picard functor of X over B
\def\Pic{\mathop{\mathrm{Pic}}\nolimits}
\def\Picardstack{\mathcal{P}\!\mathit{ic}}
\def\Picardfunctor{\mathrm{Pic}}
\def\Deformationcategory{\mathcal{D}\!\mathit{ef}}

\setlength{\emergencystretch}{3em}
\begin{document}
\title[Stacks Project, Tag 02UZ]{640 --- Grothendieck vanishing on Noetherian topological spaces}
\maketitle
\noindent Copyright (C) 2005--2025 Johan de Jong. Stacks Project authors. Source: \href{https://stacks.math.columbia.edu/tag/02UZ}{Tag 02UZ}.

\noindent Editorial difficulty note (not author text). Difficulty H2: The author reduces arbitrary sheaves to extension-by-zero constant generators, decomposes irreducible components, and inducts on dimension using the closed complement. The finite-generation filtration and colimit reduction form a substantial proof beyond ordinary qualification results..

\noindent Editorial provenance note (not author text). History: excerpt from revision \texttt{a0444\allowbreak 6e57e\allowbreak c1fbc\allowbreak 252a8\allowbreak 71afc\allowbreak ec775\allowbreak 2fb28\allowbreak 07b14}, cohomology.tex, lines 3665--3733. Reference presentation was adapted; original-author context and core support blocks were retained or added. The mathematical wording is unchanged. Licensed under GNU FDL 1.2 or later; no invariant sections or cover texts. See ../licenses/COPYING. Context blocks reproduce author definitions and supporting results; they are not additional counted problems.

\begin{definition}
\label{modules-definition-generated-by-local-sections}
Let $(X, \mathcal{O}_X)$ be a ringed space.
Let $\mathcal{F}$ be a sheaf of $\mathcal{O}_X$-modules.
Given a set $I$, and
local sections $s_i$, $i \in I$ of $\mathcal{F}$
we say that the subsheaf $\mathcal{G}$ of
Lemma \href{https://stacks.math.columbia.edu/tag/01AP}{lemma generated by local sections (Tag 01AP)}
above is the {\it subsheaf generated by the $s_i$}.
\end{definition}

\begin{lemma}
\label{lemma-vanishing-generated-one-section}
\begin{reference}
This is a special case of \href{https://stacks.math.columbia.edu/bibliography}{Bibliography: Tohoku (Proposition 3.6.1)}.
\end{reference}
Let $X$ be a topological space. Let $d \geq 0$ be an integer. Assume
\begin{enumerate}
\item $X$ is quasi-compact,
\item the quasi-compact opens form a basis for $X$, and
\item the intersection of two quasi-compact opens is quasi-compact.
\item $H^p(X, j_!\underline{\mathbf{Z}}_U) = 0$ for all $p > d$
and any quasi-compact open $j : U \to X$.
\end{enumerate}
Then $H^p(X, \mathcal{F}) = 0$ for all $p > d$
and any abelian sheaf $\mathcal{F}$ on $X$.
\end{lemma}

\begin{proof}
Let $S = \coprod_{U \subset X} \mathcal{F}(U)$ where $U$ runs over the
quasi-compact opens of $X$.
For any finite subset $A = \{s_1, \ldots, s_n\} \subset S$,
let $\mathcal{F}_A$ be the subsheaf of $\mathcal{F}$ generated
by all $s_i$ (see
Modules, Definition \ref{modules-definition-generated-by-local-sections}).
Note that if $A \subset A'$, then $\mathcal{F}_A \subset \mathcal{F}_{A'}$.
Hence $\{\mathcal{F}_A\}$ forms a system over the
directed partially ordered set of finite subsets of $S$.
By Modules, Lemma \href{https://stacks.math.columbia.edu/tag/01AR}{lemma generated by local sections stalk (Tag 01AR)}
it is clear that
$$
\colim_A \mathcal{F}_A = \mathcal{F}
$$
by looking at stalks. By
Lemma \ref{lemma-quasi-separated-cohomology-colimit} we have
$$
H^p(X, \mathcal{F}) =
\colim_A H^p(X, \mathcal{F}_A)
$$
Hence it suffices to prove the vanishing for the abelian sheaves
$\mathcal{F}_A$. In other words, it suffices to prove the
result when $\mathcal{F}$ is generated by finitely many local sections
over quasi-compact opens of $X$.

\medskip\noindent
Suppose that $\mathcal{F}$ is generated by the local sections
$s_1, \ldots, s_n$. Let $\mathcal{F}' \subset \mathcal{F}$
be the subsheaf generated by $s_1, \ldots, s_{n - 1}$.
Then we have a short exact sequence
$$
0 \to \mathcal{F}' \to \mathcal{F} \to \mathcal{F}/\mathcal{F}' \to 0
$$
From the long exact sequence of cohomology we see that it suffices
to prove the vanishing for the abelian sheaves $\mathcal{F}'$
and $\mathcal{F}/\mathcal{F}'$ which are generated by fewer than
$n$ local sections. Hence it suffices to prove the vanishing
for sheaves generated by at most one local section. These sheaves
are exactly the quotients of the sheaves $j_!\underline{\mathbf{Z}}_U$
where $U$ is a quasi-compact open of $X$.

\medskip\noindent
Assume now that we have a short exact sequence
$$
0 \to \mathcal{K} \to j_!\underline{\mathbf{Z}}_U \to \mathcal{F} \to 0
$$
with $U$ quasi-compact open in $X$.
It suffices to show that $H^q(X, \mathcal{K})$ is zero for $q \geq d + 1$.
As above we can write $\mathcal{K}$ as the filtered colimit of
subsheaves $\mathcal{K}'$ generated by finitely many sections over
quasi-compact opens. Then $\mathcal{F}$ is the filtered colimit of the
sheaves $j_!\underline{\mathbf{Z}}_U/\mathcal{K}'$. In this way we
reduce to the case that $\mathcal{K}$ is generated by finitely many
sections over quasi-compact opens. Note that $\mathcal{K}$
is a subsheaf of $\underline{\mathbf{Z}}_X$. Thus by
Lemma \ref{lemma-subsheaf-of-constant-sheaf} there exists a finite
filtration of $\mathcal{K}$ whose successive quotients $\mathcal{Q}$ fit
into a short exact sequence
$$
0 \to j''_!\underline{\mathbf{Z}}_W \to
j'_!\underline{\mathbf{Z}}_V \to \mathcal{Q} \to 0
$$
with $j'' : W \to X$ and $j' : V \to X$ the inclusions of quasi-compact opens.
Hence the vanishing of $H^p(X, \mathcal{Q})$ for $p > d$ follows
from our assumption (in the lemma) on the vanishing of the cohomology groups
of $j''_!\underline{\mathbf{Z}}_W$ and $j'_!\underline{\mathbf{Z}}_V$.
Returning to $\mathcal{K}$ this, via an induction argument using the
long exact cohomology sequence, implies the desired vanishing for it as well.
\end{proof}


\begin{lemma}
\label{lemma-subsheaf-of-constant-sheaf}
\begin{reference}
\href{https://stacks.math.columbia.edu/bibliography}{Bibliography: Tohoku (Page 168)}.
\end{reference}
Let $X$ be a topological space such that the intersection of any
two quasi-compact opens is quasi-compact. Let
$\mathcal{F} \subset \underline{\mathbf{Z}}$
be a subsheaf generated by finitely many sections over quasi-compact opens.
Then there exists a finite filtration
$$
(0) = \mathcal{F}_0 \subset \mathcal{F}_1 \subset \ldots \subset
\mathcal{F}_n = \mathcal{F}
$$
by abelian subsheaves such that for each $0 < i \leq n$
there exists a short exact sequence
$$
0 \to j'_!\underline{\mathbf{Z}}_V \to j_!\underline{\mathbf{Z}}_U \to
\mathcal{F}_i/\mathcal{F}_{i - 1} \to 0
$$
with $j : U \to X$ and $j' : V \to X$ the inclusion of quasi-compact opens
into $X$.
\end{lemma}

\begin{proof}
Say $\mathcal{F}$ is generated by the sections $s_1, \ldots, s_t$ over the
quasi-compact opens $U_1, \ldots, U_t$. Since $U_i$ is quasi-compact and
$s_i$ a locally constant function to $\mathbf{Z}$ we may assume, after
possibly replacing $U_i$ by the parts of a finite decomposition into open
and closed subsets, that $s_i$ is a constant section.
Say $s_i = n_i$ with $n_i \in \mathbf{Z}$. Of course we can remove
$(U_i, n_i)$ from the list if $n_i = 0$. Flipping signs if necessary
we may also assume $n_i > 0$. Next, for any subset $I \subset \{1, \ldots, t\}$
we may add $\bigcap_{i \in I} U_i$ and $\gcd(n_i, i \in I)$ to the list.
After doing this we see that our list $(U_1, n_1), \ldots, (U_t, n_t)$
satisfies the following property:
For $x \in X$ set $I_x = \{i \in \{1, \ldots, t\} \mid x \in U_i\}$.
Then $\gcd(n_i, i \in I_x)$ is attained by $n_i$ for some $i \in I_x$.

\medskip\noindent
As our filtration we take $\mathcal{F}_0 = (0)$ and
$\mathcal{F}_n$ generated by the sections $n_i$ over $U_i$ for those
$i$ such that $n_i \leq n$. It is clear that
$\mathcal{F}_n = \mathcal{F}$ for $n \gg 0$. Moreover, the quotient
$\mathcal{F}_n/\mathcal{F}_{n - 1}$ is generated by the section
$n$ over $U = \bigcup_{n_i \leq n} U_i$ and the kernel of the map
$j_!\underline{\mathbf{Z}}_U \to \mathcal{F}_n/\mathcal{F}_{n - 1}$
is generated by the section $n$ over $V = \bigcup_{n_i \leq n - 1} U_i$.
Thus a short exact sequence as in the statement of the lemma.
\end{proof}


\begin{lemma}
\label{lemma-quasi-separated-cohomology-colimit}
Let $X$ be a ringed space. Assume that the underlying topological space
of $X$ has the following properties:
\begin{enumerate}
\item there exists a basis of quasi-compact open subsets, and
\item the intersection of any two quasi-compact opens is quasi-compact.
\end{enumerate}
Then for any directed system $(\mathcal{F}_i, \varphi_{ii'})$
of sheaves of $\mathcal{O}_X$-modules and for any quasi-compact open
$U \subset X$ the canonical map
$$
\colim_i H^q(U, \mathcal{F}_i)
\longrightarrow
H^q(U, \colim_i \mathcal{F}_i)
$$
is an isomorphism for every $q \geq 0$.
\end{lemma}

\begin{proof}
It is important in this proof to argue for all quasi-compact opens
$U \subset X$ at the same time.
The result is true for $q = 0$ and any quasi-compact open $U \subset X$ by
Sheaves, Lemma \href{https://stacks.math.columbia.edu/tag/009F}{lemma directed colimits sections (Tag 009F)}
(combined with
Topology, Lemma \href{https://stacks.math.columbia.edu/tag/0069}{lemma topology quasi separated scheme (Tag 0069)}).
Assume that we have proved the result for all $q \leq q_0$ and let
us prove the result for $q = q_0 + 1$.

\medskip\noindent
By our conventions on directed systems the index set $I$ is directed,
and any system of $\mathcal{O}_X$-modules $(\mathcal{F}_i, \varphi_{ii'})$
over $I$ is directed.
By Injectives, Lemma \href{https://stacks.math.columbia.edu/tag/01DI}{lemma sheaves modules space (Tag 01DI)} the category
of $\mathcal{O}_X$-modules has functorial injective embeddings.
Thus for any system $(\mathcal{F}_i, \varphi_{ii'})$ there exists a
system $(\mathcal{I}_i, \varphi_{ii'})$ with each $\mathcal{I}_i$ an
injective $\mathcal{O}_X$-module and a morphism of systems given
by injective $\mathcal{O}_X$-module maps
$\mathcal{F}_i \to \mathcal{I}_i$. Denote $\mathcal{Q}_i$ the
cokernel so that we have short exact sequences
$$
0 \to
\mathcal{F}_i \to
\mathcal{I}_i \to
\mathcal{Q}_i \to 0.
$$
We claim that the sequence
$$
0 \to
\colim_i \mathcal{F}_i \to
\colim_i \mathcal{I}_i \to
\colim_i \mathcal{Q}_i \to 0.
$$
is also a short exact sequence of $\mathcal{O}_X$-modules.
We may check this on stalks. By
Sheaves, Sections \href{https://stacks.math.columbia.edu/tag/009D}{section limits presheaves (Tag 009D)}
and \href{https://stacks.math.columbia.edu/tag/009E}{section limits sheaves (Tag 009E)}
taking stalks commutes with colimits. Since a directed colimit
of short exact sequences of abelian groups is short exact
(see Algebra, Lemma \href{https://stacks.math.columbia.edu/tag/00DB}{lemma directed colimit exact (Tag 00DB)})
we deduce the result. We claim that
$H^q(U, \colim_i \mathcal{I}_i) = 0$ for all quasi-compact
open $U \subset X$ and all $q \geq 1$. Accepting this claim
for the moment consider the diagram
$$
\xymatrix{
\colim_i H^{q_0}(U, \mathcal{I}_i) \ar[d] \ar[r] &
\colim_i H^{q_0}(U, \mathcal{Q}_i) \ar[d] \ar[r] &
\colim_i H^{q_0 + 1}(U, \mathcal{F}_i) \ar[d] \ar[r] &
0 \ar[d] \\
H^{q_0}(U, \colim_i \mathcal{I}_i) \ar[r] &
H^{q_0}(U, \colim_i \mathcal{Q}_i) \ar[r] &
H^{q_0 + 1}(U, \colim_i \mathcal{F}_i) \ar[r] &
0
}
$$
The zero at the lower right corner comes from the claim and the
zero at the upper right corner comes from the fact that the sheaves
$\mathcal{I}_i$ are injective.
The top row is exact by an application of
Algebra, Lemma \href{https://stacks.math.columbia.edu/tag/00DB}{lemma directed colimit exact (Tag 00DB)}.
Hence by the snake lemma we deduce the
result for $q = q_0 + 1$.

\medskip\noindent
It remains to show that the claim is true. We will use
Lemma \ref{lemma-cech-vanish-basis}.
Let $\mathcal{B}$ be the collection of all quasi-compact open
subsets of $X$. This is a basis for the topology on $X$ by assumption.
Let $\text{Cov}$ be the collection of finite open coverings
$\mathcal{U} : U = \bigcup_{j = 1, \ldots, m} U_j$ with each
of $U$, $U_j$ quasi-compact open in $X$. By the result for $q = 0$
we see that for $\mathcal{U} \in \text{Cov}$ we have
$$
\check{\mathcal{C}}^\bullet(\mathcal{U}, \colim_i \mathcal{I}_i)
=
\colim_i \check{\mathcal{C}}^\bullet(\mathcal{U}, \mathcal{I}_i)
$$
because all the multiple intersections $U_{j_0 \ldots j_p}$
are quasi-compact. By Lemma \href{https://stacks.math.columbia.edu/tag/01EP}{lemma injective trivial cech (Tag 01EP)}
each of the complexes in the colimit of {\v C}ech complexes is
acyclic in degree $\geq 1$. Hence by
Algebra, Lemma \href{https://stacks.math.columbia.edu/tag/00DB}{lemma directed colimit exact (Tag 00DB)}
we see that also the {\v C}ech complex
$\check{\mathcal{C}}^\bullet(\mathcal{U}, \colim_i \mathcal{I}_i)$
is acyclic in degrees $\geq 1$. In other words we see that
$\check{H}^p(\mathcal{U},  \colim_i \mathcal{I}_i) = 0$
for all $p \geq 1$. Thus the assumptions of
Lemma \ref{lemma-cech-vanish-basis} are satisfied and the claim follows.
\end{proof}


\begin{lemma}
\label{lemma-cech-vanish-basis}
(Variant of Lemma \href{https://stacks.math.columbia.edu/tag/01EV}{lemma cech vanish (Tag 01EV)}.)
Let $X$ be a ringed space.
Let $\mathcal{B}$ be a basis for the topology on $X$.
Let $\mathcal{F}$ be an $\mathcal{O}_X$-module.
Assume there exists a set of open coverings $\text{Cov}$
with the following properties:
\begin{enumerate}
\item For every $\mathcal{U} \in \text{Cov}$
with $\mathcal{U} : U = \bigcup_{i \in I} U_i$ we have
$U, U_i \in \mathcal{B}$ and every $U_{i_0 \ldots i_p} \in \mathcal{B}$.
\item For every $U \in \mathcal{B}$ the open coverings of $U$
occurring in $\text{Cov}$ is a cofinal system of open coverings
of $U$.
\item For every $\mathcal{U} \in \text{Cov}$ we have
$\check{H}^p(\mathcal{U}, \mathcal{F}) = 0$ for all $p > 0$.
\end{enumerate}
Then $H^p(U, \mathcal{F}) = 0$ for all $p > 0$ and any $U \in \mathcal{B}$.
\end{lemma}

\begin{proof}
Let $\mathcal{F}$ and $\text{Cov}$ be as in the lemma.
We will indicate this by saying ``$\mathcal{F}$ has vanishing higher
{\v C}ech cohomology for any $\mathcal{U} \in \text{Cov}$''.
Choose an embedding $\mathcal{F} \to \mathcal{I}$ into an
injective $\mathcal{O}_X$-module.
By Lemma \href{https://stacks.math.columbia.edu/tag/01EP}{lemma injective trivial cech (Tag 01EP)} $\mathcal{I}$
has vanishing higher {\v C}ech cohomology for any $\mathcal{U} \in \text{Cov}$.
Let $\mathcal{Q} = \mathcal{I}/\mathcal{F}$
so that we have a short exact sequence
$$
0 \to \mathcal{F} \to \mathcal{I} \to \mathcal{Q} \to 0.
$$
By Lemma \href{https://stacks.math.columbia.edu/tag/01EU}{lemma ses cech h1 (Tag 01EU)} and our assumption (2)
this sequence gives rise to an exact sequence
$$
0 \to \mathcal{F}(U) \to \mathcal{I}(U) \to \mathcal{Q}(U) \to 0.
$$
for every $U \in \mathcal{B}$. Hence for any $\mathcal{U} \in \text{Cov}$
we get a short exact sequence of {\v C}ech complexes
$$
0 \to
\check{\mathcal{C}}^\bullet(\mathcal{U}, \mathcal{F}) \to
\check{\mathcal{C}}^\bullet(\mathcal{U}, \mathcal{I}) \to
\check{\mathcal{C}}^\bullet(\mathcal{U}, \mathcal{Q}) \to 0
$$
since each term in the {\v C}ech complex is made up out of a product of
values over elements of $\mathcal{B}$ by assumption (1).
In particular we have a long exact sequence of {\v C}ech cohomology
groups for any open covering $\mathcal{U} \in \text{Cov}$.
This implies that $\mathcal{Q}$ is also an $\mathcal{O}_X$-module
with vanishing higher {\v C}ech cohomology for all
$\mathcal{U} \in \text{Cov}$.

\medskip\noindent
Next, we look at the long exact cohomology sequence
$$
\xymatrix{
0 \ar[r] &
H^0(U, \mathcal{F}) \ar[r] &
H^0(U, \mathcal{I}) \ar[r] &
H^0(U, \mathcal{Q}) \ar[lld] \\
&
H^1(U, \mathcal{F}) \ar[r] &
H^1(U, \mathcal{I}) \ar[r] &
H^1(U, \mathcal{Q}) \ar[lld] \\
&
\ldots & \ldots & \ldots \\
}
$$
for any $U \in \mathcal{B}$. Since $\mathcal{I}$ is injective we
have $H^n(U, \mathcal{I}) = 0$ for $n > 0$ (see
Derived Categories, Lemma \href{https://stacks.math.columbia.edu/tag/015B}{lemma higher derived functors (Tag 015B)}).
By the above we see that $H^0(U, \mathcal{I}) \to H^0(U, \mathcal{Q})$
is surjective and hence $H^1(U, \mathcal{F}) = 0$.
Since $\mathcal{F}$ was an arbitrary $\mathcal{O}_X$-module with
vanishing higher {\v C}ech cohomology for all $\mathcal{U} \in \text{Cov}$
we conclude that also $H^1(U, \mathcal{Q}) = 0$ since $\mathcal{Q}$ is
another of these sheaves (see above). By the long exact sequence this in
turn implies that $H^2(U, \mathcal{F}) = 0$. And so on and so forth.
\end{proof}


\begin{proposition}[Grothendieck]
\label{proposition-vanishing-Noetherian}
\begin{reference}
\href{https://stacks.math.columbia.edu/bibliography}{Bibliography: Tohoku (Theorem 3.6.5)}.
\end{reference}
Let $X$ be a Noetherian topological space.
If $\dim(X) \leq d$, then $H^p(X, \mathcal{F}) = 0$
for all $p > d$ and any abelian sheaf $\mathcal{F}$
on $X$.
\end{proposition}

\begin{proof}
We prove this lemma by induction on $d$.
So fix $d$ and assume the lemma holds for all
Noetherian topological spaces of dimension $< d$.

\medskip\noindent
Let $\mathcal{F}$ be an abelian sheaf on $X$.
Suppose $U \subset X$ is an open. Let $Z \subset X$
denote the closed complement.
Denote $j : U \to X$ and $i : Z \to X$ the inclusion maps.
Then there is a short exact sequence
$$
0 \to j_{!}j^*\mathcal{F} \to \mathcal{F} \to i_*i^*\mathcal{F} \to 0
$$
see Modules, Lemma \href{https://stacks.math.columbia.edu/tag/02UT}{lemma canonical exact sequence (Tag 02UT)}.
Note that $j_!j^*\mathcal{F}$ is supported on
the topological closure $Z'$ of $U$, i.e., it is of
the form $i'_*\mathcal{F}'$ for some abelian sheaf $\mathcal{F}'$
on $Z'$, where $i' : Z' \to X$ is the inclusion.

\medskip\noindent
We can use this to reduce to the case where $X$ is irreducible.
Namely, according to
Topology, Lemma \href{https://stacks.math.columbia.edu/tag/0052}{lemma Noetherian (Tag 0052)}
$X$ has finitely
many irreducible components. If $X$ has more than one irreducible
component, then let $Z \subset X$ be an irreducible component of $X$
and set $U = X \setminus Z$. By the above, and the long exact sequence
of cohomology, it suffices to prove the vanishing of
$H^p(X, i_*i^*\mathcal{F})$ and $H^p(X, i'_*\mathcal{F}')$ for $p > d$.
By Lemma \href{https://stacks.math.columbia.edu/tag/02UV}{lemma cohomology and closed immersions (Tag 02UV)} it suffices to prove
$H^p(Z, i^*\mathcal{F})$ and $H^p(Z', \mathcal{F}')$ vanish for $p > d$.
Since $Z'$ and $Z$ have fewer irreducible components we indeed
reduce to the case of an irreducible $X$.

\medskip\noindent
If $d = 0$ and $X$ is irreducible, then $X$ is the only nonempty
open subset of $X$. Hence every sheaf is constant and higher cohomology
groups vanish (for example by
Lemma \href{https://stacks.math.columbia.edu/tag/02UW}{lemma irreducible constant cohomology zero (Tag 02UW)}).

\medskip\noindent
Suppose $X$ is irreducible of dimension $d > 0$.
By Lemma \ref{lemma-vanishing-generated-one-section}
we reduce to the case where
$\mathcal{F} = j_!\underline{\mathbf{Z}}_U$ for some open $U \subset X$.
In this case we look at the short exact sequence
$$
0 \to j_!(\underline{\mathbf{Z}}_U) \to
\underline{\mathbf{Z}}_X \to i_*\underline{\mathbf{Z}}_Z \to 0
$$
where $Z = X \setminus U$.
By Lemma \href{https://stacks.math.columbia.edu/tag/02UW}{lemma irreducible constant cohomology zero (Tag 02UW)}
we have the vanishing of $H^p(X, \underline{\mathbf{Z}}_X)$
for all $p \geq 1$. By induction we have
$H^p(X, i_*\underline{\mathbf{Z}}_Z) = H^p(Z, \underline{\mathbf{Z}}_Z) = 0$
for $p \geq d$. Hence we win by the long exact cohomology sequence.
\end{proof}
\end{document}
